% GENERATED FILE. Edit canonical lessons, then run npm run generate:books.
% canonical-source-hash: fnv1a64:1d05e52f1d2f7fad
% canonical-lessons: ML-C181-kitakkuka, ML-C181-tiriyuka, ML-C181-katakkuka, ML-C181-valikkuka, ML-C181-talluka

\chapter{To lie down, To turn, To cross, To pull, To push}
\label{ch:ml-to-lie-down-to-turn-to-cross-to-pull-to-push}
\hlchaptermodality{181}

\begin{chapteropening}
\textbf{By the end of this chapter:} \emph{I can say to lie down, to turn, to cross, to pull and to push.}

Complete the last lesson of chapter 181: I can say to lie down, to turn, to cross, to pull and to push.
\end{chapteropening}

\section[\texorpdfstring{\ml{കിടക്കുക} (kiṭakkuka)}{കിടക്കുക (kiṭakkuka)}]{\ml{കിടക്കുക} (kiṭakkuka) — to lie down}

\label{lesson:ML-C181-kitakkuka}



\begin{quote}
Before the new one: say the Malayalam for an advertisement, then the Malayalam for a receipt.
\end{quote}

\subsection*{You'll want to know: \ml{കിടക്കുക}}
\textbf{\ml{കിടക്കുക}} — \emph{kiṭakkuka} — \textquotedblleft{}to lie down\textquotedblright{}.

\begin{grammarlens}[title={What you've built}]
One more word for this chapter.
\end{grammarlens}

\subsection*{Your turn}
\begin{itemize}
\raggedright
  \item \emph{Say it:} \emph{kiṭakkuka}
  \item \emph{Say it:} \emph{kiṭakkuka}, once more
  \item \emph{Say it:} say \emph{kiṭakkuka}
  \item \emph{Recall:} say the Malayalam for an advertisement, then the Malayalam for a receipt, then say \emph{kiṭakkuka} again
\end{itemize}

\begin{tcolorbox}[breakable,colback=teal!4,colframe=teal!35!black,title={Before you move on}]
What does \ml{കിടക്കുക} mean? (\textquotedblleft{}To lie down\textquotedblright{}.) Say it once more.
\end{tcolorbox}
\section[\texorpdfstring{\ml{തിരിയുക} (tiriyuka)}{തിരിയുക (tiriyuka)}]{\ml{തിരിയുക} (tiriyuka) — to turn}

\label{lesson:ML-C181-tiriyuka}



\begin{quote}
Before the new one: say the Malayalam for a receipt, then the Malayalam for to lie down.
\end{quote}

\subsection*{You'll want to know: \ml{തിരിയുക}}
\textbf{\ml{തിരിയുക}} — \emph{tiriyuka} — \textquotedblleft{}to turn\textquotedblright{}.

\begin{grammarlens}[title={What you've built}]
One more word for this chapter.
\end{grammarlens}

\subsection*{Your turn}
\begin{itemize}
\raggedright
  \item \emph{Say it:} \emph{tiriyuka}
  \item \emph{Say it:} \emph{tiriyuka}, once more
  \item \emph{Say it:} say \emph{tiriyuka}
  \item \emph{Recall:} say the Malayalam for a receipt, then the Malayalam for to lie down, then say \emph{tiriyuka} again
\end{itemize}

\begin{tcolorbox}[breakable,colback=teal!4,colframe=teal!35!black,title={Before you move on}]
What does \ml{തിരിയുക} mean? (\textquotedblleft{}To turn\textquotedblright{}.) Say it once more.
\end{tcolorbox}
\section[\texorpdfstring{\ml{കടക്കുക} (kaṭakkuka)}{കടക്കുക (kaṭakkuka)}]{\ml{കടക്കുക} (kaṭakkuka) — to cross}

\label{lesson:ML-C181-katakkuka}



\begin{quote}
Before the new one: say the Malayalam for to lie down, then the Malayalam for to turn.
\end{quote}

\subsection*{You'll want to know: \ml{കടക്കുക}}
\textbf{\ml{കടക്കുക}} — \emph{kaṭakkuka} — \textquotedblleft{}to cross\textquotedblright{}.

\begin{grammarlens}[title={What you've built}]
One more word for this chapter.
\end{grammarlens}

\subsection*{Your turn}
\begin{itemize}
\raggedright
  \item \emph{Say it:} \emph{kaṭakkuka}
  \item \emph{Say it:} \emph{kaṭakkuka}, once more
  \item \emph{Say it:} say \emph{kaṭakkuka}
  \item \emph{Recall:} say the Malayalam for to lie down, then the Malayalam for to turn, then say \emph{kaṭakkuka} again
\end{itemize}

\begin{tcolorbox}[breakable,colback=teal!4,colframe=teal!35!black,title={Before you move on}]
What does \ml{കടക്കുക} mean? (\textquotedblleft{}To cross\textquotedblright{}.) Say it once more.
\end{tcolorbox}
\section[\texorpdfstring{\ml{വലിക്കുക} (valikkuka)}{വലിക്കുക (valikkuka)}]{\ml{വലിക്കുക} (valikkuka) — to pull}

\label{lesson:ML-C181-valikkuka}



\begin{quote}
Before the new one: say the Malayalam for to turn, then the Malayalam for to cross.
\end{quote}

\subsection*{You'll want to know: \ml{വലിക്കുക}}
\textbf{\ml{വലിക്കുക}} — \emph{valikkuka} — \textquotedblleft{}to pull\textquotedblright{}.

\begin{grammarlens}[title={What you've built}]
One more word for this chapter.
\end{grammarlens}

\subsection*{Your turn}
\begin{itemize}
\raggedright
  \item \emph{Say it:} \emph{valikkuka}
  \item \emph{Say it:} \emph{valikkuka}, once more
  \item \emph{Say it:} say \emph{valikkuka}
  \item \emph{Recall:} say the Malayalam for to turn, then the Malayalam for to cross, then say \emph{valikkuka} again
\end{itemize}

\begin{tcolorbox}[breakable,colback=teal!4,colframe=teal!35!black,title={Before you move on}]
What does \ml{വലിക്കുക} mean? (\textquotedblleft{}To pull\textquotedblright{}.) Say it once more.
\end{tcolorbox}
\section[\texorpdfstring{\ml{തള്ളുക} (taḷḷuka)}{തള്ളുക (taḷḷuka)}]{\ml{തള്ളുക} (taḷḷuka) — to push}

\label{lesson:ML-C181-talluka}



\begin{quote}
Before the new one: say the Malayalam for to cross, then the Malayalam for to pull.
\end{quote}

\subsection*{You'll want to know: \ml{തള്ളുക}}
\textbf{\ml{തള്ളുക}} — \emph{taḷḷuka} — \textquotedblleft{}to push\textquotedblright{}.

\begin{grammarlens}[title={What you've built}]
One more word for this chapter.
\end{grammarlens}

\subsection*{Your turn}
\begin{itemize}
\raggedright
  \item \emph{Say it:} \emph{taḷḷuka}
  \item \emph{Say it:} \emph{taḷḷuka}, once more
  \item \emph{Say it:} say \emph{taḷḷuka}
  \item \emph{Recall:} say the Malayalam for to cross, then the Malayalam for to pull, then say \emph{taḷḷuka} again
\end{itemize}

\begin{tcolorbox}[breakable,colback=teal!4,colframe=teal!35!black,title={Before you move on}]
What does \ml{തള്ളുക} mean? (\textquotedblleft{}To push\textquotedblright{}.) Say it once more.
\end{tcolorbox}
